\documentclass[12pt]{article}
\usepackage[utf8]{inputenc}
\usepackage[T1]{fontenc}
\usepackage{hyperref}
\title{Unggaranggu Language Resources for K-12 Curriculum Development}
\author{Compiled by Claude Sonnet 5.5 (claude-sonnet-5-5) for Andr\'as Kornai}
\date{}
\begin{document}
\maketitle
\section{Introduction}
Unggaranggu (Unggarrangu, Ongkarango; ISO 639-3 xun) is an extinct Worrorran language of the Kimberley region of Western Australia. Very little is known of it; what exists was recorded by Howard Coate in the 1960s. The Unggarrangu people and the Yawijibaya are the traditional owners of the Buccaneer Archipelago, known together as the Mayala group. Nothing educational was found in the search results.
\section{Primary Educational Resources}
None found.
\section{Linguistic and Reference Materials}
\begin{itemize}
\item Wikipedia, Unggarrangu: \url{https://en.wikipedia.org/wiki/Unggarrangu}
\item Glottolog entry (Western Worrorran, extinct): \url{https://glottolog.org/resource/languoid/id/ungg1243}
\end{itemize}
\section{Religious and Community Resources}
No Unggaranggu-specific organisation was found. The neighbouring but distinct Wunambal Gaambera language is represented by the Wunambal Gaambera Aboriginal Corporation; its materials (including a 2023 dictionary) must not be treated as Unggaranggu. See AIATSIS Austlang for Gaambera: \url{https://collection.aiatsis.gov.au/austlang/language/K39}
\section{Digital and Audio Resources}
None found.
\section{Curriculum Development Context}
Material for a language curriculum does not exist in the sources found. Any school work would be cultural and historical (the Mayala people, the Buccaneer Archipelago) and would need to be led by Mayala traditional owners. Archival sources such as Coate's notes would have to be sought through AIATSIS or similar archives, which was not verified here.
\end{document}